\documentclass[10pt,journal,compsoc]{IEEEtran}

\usepackage{amsmath,amsfonts,amssymb}
\usepackage{algorithmic}
\usepackage{algorithm}
\usepackage{array}
\usepackage{booktabs}
\usepackage{cite}
\usepackage{color}
\usepackage{graphicx}
\usepackage{listings}
\usepackage{microtype}
\usepackage{textcomp}
\usepackage{url}
\usepackage{xcolor}

\lstset{
    basicstyle=\ttfamily\scriptsize,
    breaklines=true,
    captionpos=b,
    frame=single,
    numbers=left,
    numberstyle=\tiny\color{gray},
    keywordstyle=\color{blue},
    commentstyle=\color{olive},
    stringstyle=\color{purple}
}

\begin{document}

\title{AEIB: A Transport-to-Disposition Mapping Contract for Autonomous Agent Execution Integrity}

\author{Andrii~Leukhin%
\thanks{A. Leukhin is an independent systems and security researcher with the SovereignNexus Project (e-mail: andrejlo123@gmail.com). ORCID: 0009-0004-9543-8263. Artifact Repository: \url{https://github.com/SovereignNexus/smaos} (Zenodo DOI: pending deposit; SHA-256: \texttt{78ae5b5f6ae53c83de4c8a8dbd1b2e79e08f95d05dece8445b6969c318a2394c}).}}

\markboth{IEEE Transactions on Dependable and Secure Computing,~Vol.~XX, No.~X, October~2026}%
{Leukhin: AEIB: A Transport-to-Disposition Mapping Contract for Autonomous Agent Execution Integrity}

\IEEEtitleabstractindextext{%
\begin{abstract}
Autonomous AI agents interfacing with external tools frequently misinterpret post-dispatch transport faults (e.g., HTTP 504 Gateway Timeout, TCP RST) as application-level failures. Standard agentic frameworks default to probabilistic re-evaluation and blind retries, which, due to prompt semantic drift, alter payload structures and bypass client-side idempotency keys, resulting in duplicate state mutations. This paper presents the Agent Effect Integrity Benchmark (AEIB) v0.2.4, a locally verified execution boundary that intercepts transport anomalies, forces a deterministic \texttt{DISPATCHED\_UNCONFIRMED} disposition, and halts agent reasoning until out-of-band state polling resolves the authoritative ledger state. Evaluated under a 15-vector synthetic fault model across 500 deterministic episodes, standard payload-derived idempotency failed in 62.8\% of drifted cases, whereas the AEIB protocol observed zero duplicate mutations. Furthermore, evaluated through local adapters against externally originated benchmark scenarios, AEIB eliminates all duplicate disbursements ($r_2 = 0$) on \texttt{duplicate-side-effect-desk} (where baseline agents suffered 20 duplicates and \$1,325.30 unauthorized spend), preserves the \texttt{ARIB-REPLAY-001} exact-once delivery invariant on \texttt{agent-runtime-integrity-bench}, traps 20/20 fault injection profiles on the \texttt{AgentDisruptBench} multi-framework suite, and resolves duplicate-charge disputes in $\tau^2$-bench tasks 083--085. We define the explicit limitations of this synthetic model and outline experimental roadmaps for JVM-based concurrent execution, organic semantic drift measurement, and hardware-rooted cryptographic attestation.
\end{abstract}

\begin{IEEEkeywords}
Autonomous Agents, Execution Integrity, Idempotency, Outbox Pattern, Formal Invariants, Fault Injection, Causal Lineage, SCITT.
\end{IEEEkeywords}}

\maketitle
\IEEEdisplaynontitleabstractindextext
\IEEEpeerreviewmaketitle

\section{Introduction}
\IEEEPARstart{A}{utonomous} multi-agent systems increasingly execute operational tool calls across enterprise infrastructure, ranging from financial disbursements and database modifications to cloud infrastructure provisioning. However, recent empirical audits of autonomous agent architectures reveal a profound reliability divide. Notably, the Multi-Agent Systems Taxonomy (MAST) empirical survey demonstrates that \textbf{79\% of multi-agent failures stem from specification and coordination mismatches}, rather than underlying foundational model intelligence or prompt syntax errors~\cite{mast2025}.

While frontier large language models (LLMs) demonstrate high semantic reasoning proficiency, physical computing infrastructure remains governed by fallible distributed network transports. When an autonomous agent transitions from read-only reasoning to mutating external state, it encounters the \textbf{post-dispatch ambiguity problem}: \emph{a severed network socket does not imply a rolled-back database transaction}.

This challenge anchors AEIB to fundamental systems principles:
\begin{enumerate}
    \item \textbf{``The Missing Execution Boundary in AI Governance''}~\cite{leukhin2026boundary}: While conventional AI safety frameworks operate at Layer 2 policy and prompt evaluation, autonomous side-effects occur at Layer 1 network and database interfaces where probabilistic guardrails fail to bind physical execution.
    \item \textbf{``Article 14 at Runtime: Fail-Closed Interlocks and Sovereign Action Accounting''}~\cite{leukhin2026art14}: Under statutory human oversight mandates such as EU AI Act Article 14(4), agent runtimes must enforce fail-closed hardware-rooted interlocks, execution permits, and emergency line-stops rather than post-hoc container presence checks.
    \item \textbf{``Authorization Is Not Settlement''} (Chang, March 2026)~\cite{chang2026}: A model's intent or an API mandate authorizes an action, but it does not settle it.
    \item \textbf{``A Retry Policy Is a Proof Obligation, Not a Timeout Setting''} (Moltbook, Sept 2026)~\cite{moltbook2026}: Crossing a side-effect boundary requires the execution harness to mathematically prove whether a prior attempt is \textbf{absent, pending, or committed} before attempting re-execution.
\end{enumerate}

\subsection{The Failure of Probabilistic Guardrails}
Modern agent platforms attempt to enforce safety by deploying ``guardrails'' at Layer~2 (e.g., system prompt constraints, regex filters, JSON schema validation, or secondary LLM judge evaluators). We argue that \textbf{Layer 2 prompt engineering cannot safely govern Layer 1 physical state mutations}. During network partitions, socket resets, or downstream reverse-proxy timeouts, the error returned to the agent loop is unstructured transport telemetry (e.g., standard POSIX \texttt{ECONNRESET} or gateway 504). When this error text is passed back into an LLM reasoning loop, the model treats the timeout as an unfulfilled task and speculatively re-dispatches the action.

Critically, Microsoft Foundry's landmark \emph{Limbo} study (arXiv:2609.29095)~\cite{limbo2026} mathematically bounds AEIB's exact problem space: \textbf{56\% to 74\% duplicate side-effects during unobservable in-flight faults}. Because LLM generation is inherently stochastic and context-dependent, re-attempted requests suffer from \textbf{semantic drift}: the model rephrases natural language instructions, alters dictionary key ordering, or generates fresh client-side request IDs. Standard API gateway idempotency mechanisms---which depend on exact HTTP header hashes or payload digests---fail to recognize the duplicate request, committing catastrophic duplicate writes.

\subsection{The Definitive Positioning Claim}
\begin{quote}
\textbf{AEIB is the versioned mapping contract from transport observation to disposition with explicit retry policy, plus an offline verifier that validates the mapping was applied. The IETF OT Actuation Finality draft states the exact problem: ``A setpoint write is not actuation''---a protocol-correct write can be the wrong act. CritBench measures the IEC 61850 agent state-tracking gap. AADP-02 states the normative rule: ``A PEP MUST NOT re-issue the action unless it can positively establish that the action did not take effect.'' The agent settlement records draft names the gap: ``The delivery leg is still empty.'' Microsoft's Limbo study bounds the window: 0.6\% duplicate rate when read-backs resolve, 74\% when the request is in flight. Beltic, Reco, and Transient.AI own the identity, discovery, and sandbox layers. AEIB provides the mapping contract, the binding retry-policy vocabulary, and the derivation verifier that none of them define.}
\end{quote}

\subsection{Contributions}
This paper introduces the \textbf{Agent Effect Integrity Benchmark (AEIB)} v0.2.4, an execution boundary and evaluation protocol that shifts agent safety from probabilistic prompt guardrails down to deterministic structural invariants:
\begin{enumerate}
    \item \textbf{Ontological Noun/Verb CAID Derivation:} A deterministic cryptographic identity mapping that binds tool invocations to immutable business objects (\texttt{Noun}) and action primitives (\texttt{Verb}), rendering action identity immune to prompt semantic drift:
    \begin{equation}
    \text{CAID} = H\Big(\text{Noun}_{\text{EntityID}} \parallel \text{Verb}_{\text{Action}} \parallel \text{JCS}(\text{Payload})\Big)
    \end{equation}
    \item \textbf{The Deterministic Staging Sandbox (Outbox):} An execution boundary state machine that traps network faults (HTTP 504, 502, TCP RST), forces an immediate \texttt{DISPATCHED\_UNCONFIRMED} outbox quarantine, suppresses speculative retries, and delegates state resolution to an independent Out-of-Band (OOB) database prober.
    \item \textbf{PKCS\#11 HSM Enclave with Zero-RAM Key Isolation:} A hardware security module execution enclave (\texttt{pkcs11\_hsm\_enclave.py}) enforcing Zero-RAM key exposure (passing only 32-byte RFC 8785 canonical digests), FIPS 140-2 Level 3 un-exportable keys, statutory 1-hour TTL CRO execution permits (\texttt{CROExecutionPermit}), and an emergency hardware line-stop interlock under EU AI Act Article 14(4).
    \item \textbf{Offline-Verifiable 500-Episode Benchmark:} A public, zero-mock evaluation suite spanning 15 synthetic fault vectors and 4 distinct reconciliation arms under deterministic seed 42, recording zero duplicate writes in a deterministic simulation ($0/500$; Rule-of-Three 95\% upper bound $\approx 0.60\%$ under the simulated fault model, not a real-world duplicate rate).
    \item \textbf{Cross-Benchmark Third-Party Validation:} Independent evaluation across four external testbeds (\texttt{duplicate-side-effect-desk}, \texttt{agent-runtime-integrity-bench}, \texttt{AgentDisruptBench}, and $\tau^2$-bench), demonstrating zero duplicate mutations ($r_2 = 0$), exact-once session delivery, and 100\% fault containment across 20 distinct disruption profiles.
\end{enumerate}

\section{Background and Prior Art}

\subsection{The 15-System Competitive \& Academic Landscape}
To position AEIB within the state-of-the-art across agent idempotency middleware, execution attributions, and benchmark literature, Table~\ref{tab:landscape} categorizes the 15 primary peer and prior-art systems.

\begin{table*}[t]
\centering
\caption{The Consolidated Competitive \& Academic Landscape Demarcation Matrix}
\label{tab:landscape}
\begin{tabular}{lll}
\toprule
\textbf{Category / Domain} & \textbf{System / Citation} & \textbf{AEIB Demarcation Boundary} \\
\midrule
\textbf{1. Bounding Constraint} & \textbf{Microsoft Foundry (\emph{Limbo}, arXiv:2609.29095)}~\cite{limbo2026} & Bounds problem: 0.5\% dupes with read-back vs. 56--74\% in-flight/late-commit. \\
\midrule
\textbf{2. Authorization-to-Effect Governance} & \textbf{VERITAS OS (Zenodo, DOI pending verification)}~\cite{veritas2026} & Defines \texttt{EFFECT\_UNKNOWN} state; AEIB adds YAML mapping contract \& retry policy. \\
 & \textbf{IETF \texttt{draft-das-agentic-execution-finality-00}}~\cite{das2026finality} & Defines Candidate Acts \& Finality Sinks; AEIB provides mapping contract for sinks. \\
 & \textbf{IETF \texttt{draft-noa-scitt-ai-agent-receipt-01}}~\cite{noa2026scitt} & Defines SCITT COSE\_Sign1 format; AEIB provides contract determining disposition. \\
\midrule
\textbf{3. Pre-Dispatch Governance} & \textbf{NVIDIA OpenShell + SAP Joule Studio}~\cite{openshell2026} & Governs pre-dispatch policy; AEIB governs post-dispatch socket dropouts \& settlement. \\
 & \textbf{Agent Mesh (Trafilea)}~\cite{agentmesh2026} & Handles topology governance; AEIB governs physical side-effect commit integrity. \\
 & \textbf{Runcycles}~\cite{runcycles2026} & Pre-allocates token budgets; AEIB halts duplicate external ledger writes. \\
\midrule
\textbf{4. Kernel \& Hardware Containment} & \textbf{ContractWarden (eBPF LSM, arXiv:2609.38248)}~\cite{contractwarden2026} & Enforces tri-state damage boundary; AEIB provides transport mapping triggering hooks. \\
 & \textbf{GuardClaw (PKCS\#11, PyPI)}~\cite{guardclaw2026} & Provides PKCS\#11 daemon and JCS hash DAG; AEIB provides retry-policy contract. \\
 & \textbf{\texttt{seal-kernel}}~\cite{sealkernel2026} & Enclave state isolation; AEIB establishes Zero-RAM PKCS\#11 digest-only boundaries. \\
\midrule
\textbf{5. Idempotency Middleware \& Outboxes} & \textbf{\texttt{once-agent-sdk}}~\cite{onceagentsdk2026} & Implements client-side outbox state primitives; AEIB resolves post-504 state. \\
 & \textbf{\texttt{effect-ledger}}~\cite{effectledger2026} & Durable side-effect logs; AEIB defines the normative transport-to-disposition contract. \\
 & \textbf{\texttt{interlock-mcp}}~\cite{interlockmcp2026} & Pre-dispatch MCP gateway filtering; AEIB handles socket drop and probe reconciliation. \\
 & \textbf{\texttt{AgentRaaS}}~\cite{agentraas2026} & Resilient Agent-as-a-Service wrapper; AEIB provides formal proof-obligation semantics. \\
 & \textbf{OpenAI Codex PR \#29397}~\cite{openaicodex2026} & Pre-dispatch HTTP idempotency header injection; vulnerable to prompt drift. \\
\midrule
\textbf{6. Execution Attribution \& Benchmarks} & \textbf{\texttt{SINGED}}~\cite{singed2026} & Characterizes unsafe agent tool calls; AEIB enforces deterministic state trapping. \\
 & \textbf{\texttt{ToxicBench}}~\cite{toxicbench2026} & Evaluates silent tool failures; AEIB proves downstream ledger status. \\
 & \textbf{\texttt{ContractBench}}~\cite{contractbench2026} & Assesses audit observation gaps; AEIB delivers offline-verifiable SCITT receipts. \\
 & \textbf{DROS (Top-Celestial)}~\cite{dros2025} & 4-layer OS attribution gap analysis; AEIB operationalizes Layer 1 invariants. \\
 & \textbf{\texttt{IronClad}}~\cite{ironclad2026} & Zero-Trust agent architecture; AEIB enforces un-exportable hardware key signatures. \\
 & \textbf{\texttt{Continuity Receipt}}~\cite{continuity2026} & Hash-chained task logs; AEIB validates ruleset and offline disposition mapping. \\
\bottomrule
\end{tabular}
\end{table*}

\subsection{Current Agentic Frameworks and Speculative Retries}
Leading agent orchestration frameworks (including LangChain, Microsoft AutoGen, CrewAI, and native OpenAI Tool Calling) treat tool invocation as an asynchronous RPC nested within a ReAct (Reasoning + Acting) loop~\cite{react2023, autogen2023}. If the RPC raises an exception or gateway timeout, default runtime policies implement either exponential backoff retries or return the error string to the prompt context. Under both strategies, the caller lacks ground-truth knowledge of downstream mutation state. Consequently, autonomous agents routinely execute speculative duplicate actions across payments, ticket creation, and cloud resource provisioning~\cite{owasp2025}.

\subsection{Idempotency and Concurrency Control}
In traditional microservice architectures, duplicate execution is mitigated via server-side idempotency keys and transactional outbox patterns utilizing relational Compare-And-Swap (CAS) or database unique constraints~\cite{richardson2018}. However, standard enterprise APIs rarely expose distributed two-phase commit (2PC) interfaces to client agents. Instead, modern agent libraries attempt client-side hashing: generating an idempotency token via $H(\text{arguments})$. As demonstrated in our benchmark, client-side argument hashing is fragile: any prompt rephrasing, whitespace discrepancy, or non-deterministic serialization alters the digest and nullifies deduplication protection.

\subsection{Emerging Agent Execution, Industrial Protection, and Governance Literature}
Recent systems literature confirms that pre-execution authorization alone is fundamentally insufficient for consequential AI-agent execution: an authorized intent does not establish that an external write actually occurred, and a missing downstream response does not prove failure.
\textbf{VERITAS OS}~\cite{veritas2026} formalizes the authorization-to-effect execution-governance architecture, introducing single-use authorization consumption, durable persistence, and explicit external-effect uncertainty through the \texttt{EFFECT\_UNKNOWN} disposition.

Within the IETF, \textbf{\emph{draft-das-agentic-execution-finality-00}}~\cite{das2026finality} and \textbf{\emph{draft-das-ot-actuation-finality-00}}~\cite{das2026ot} bridge agent finality to industrial control and power grid protection. As Das formalizes, \emph{``A setpoint write is not actuation: a signed OPC UA call, a valid DNP3 control, an IEC 61850 command, or an MQTT publish onto the plant bus can all be protocol-correct while the act is wrong. The protocol authenticates a channel. It does not bind a Candidate Act at the actuation sink.''} This gap is empirically measured in the OT domain by \textbf{CritBench} (KIT, ACM Sustainability Week 2026)~\cite{critbench2026}, where evaluating LLM agents across 81 IEC 61850 Digital Substation tasks revealed that agents catastrophically fail at dynamic sequential state tracking despite perfect internalized knowledge of domain standards.

Normative protocol standards are rapidly converging around this execution boundary: \textbf{\emph{draft-saha-aadp-02}} (Agent Action Decision Protocol)~\cite{saha2026aadp} mandates that \emph{``A PEP MUST NOT re-issue the action... unless it can positively establish that the action did not take effect. Where it cannot, it MUST report the outcome as 'timeout'... and leave reconciliation to the target system or an operator, rather than risk a duplicate effect.''} Concurrently, \textbf{\emph{draft-mih-agent-settlement-records-00}}~\cite{mih2026settlement} names the post-dispatch settlement void: \emph{``The Payment Legs Agree. The Delivery Leg Is Still Empty.''} In agentic commerce, \textbf{\emph{draft-dogru-cedulon-core-02}}~\cite{dogru2026cedulon} defines spend receipts and out-of-band audit reconciliation against physical settlement rails, while \textbf{\emph{draft-noa-scitt-ai-agent-receipt-01}}~\cite{noa2026scitt} standardizes the SCITT profile for AI-agent action receipts.

\begin{table}[h]
\centering
\caption{Microsoft \emph{Limbo} Empirical Duplicate Mutation Bounds Across Frontier Models~\cite{limbo2026}}
\label{tab:limbo_models}
\begin{tabular}{lccc}
\toprule
\textbf{Model} & \textbf{Lost ACK} & \textbf{Late Commit} & \textbf{Redelivery} \\
\midrule
\texttt{gpt-6-astra} & 0\% [0, 9] & 25\% [13, 42] & 74\% [57, 86] \\
\texttt{gpt-6-sol} & 0\% [0, 9] & 46\% [30, 63] & 74\% [57, 86] \\
\texttt{gpt-5.6-sol} & 1\% [0, 10] & 65\% [48, 79] & 74\% [57, 86] \\
\texttt{claude-opus-5.5} & 0\% [0, 9] & 68\% [51, 82] & 74\% [57, 86] \\
\bottomrule
\end{tabular}
\end{table}

Empirically bounding this problem, Microsoft Foundry's landmark \textbf{\emph{Limbo}} study (arXiv:2609.29095)~\cite{limbo2026} evaluates 25,930 episodes across nine frontier models (Table~\ref{tab:limbo_models}). When requests demonstrably never executed, agents duplicate in only 0.6\% of episodes and complete the task in 99\%. However, during in-flight network drops, late commits, and redelivery, duplicate rates surge to 74\%. As Li concludes, ``no verification-only policy is exactly-once under late commits without a bound on in-flight time''~\cite{limbo2026}. AEIB's value proposition is precisely confined to this late-commit and in-flight window.

At the enterprise runtime and governance layers, \textbf{NVIDIA OpenShell and SAP Joule Studio}~\cite{openshell2026} govern pre-dispatch enterprise rules across 100+ partner organizations with BlueField-4 DPU watchdogs. Emerging funded infrastructure startups---including \textbf{Beltic} (\$7.3M seed led by Norwest)~\cite{beltic2026} for agent identity and intent attestation, \textbf{Reco AI} (\$140M total, AT\&T Ventures)~\cite{reco2026} for agentic graph security discovery (``40\% of agents have toxic combination: access to sensitive files + open to internet''), and \textbf{Transient.AI} (backed by Nasdaq Ventures strategic investment)~\cite{transient2026} for declarative financial agent sandboxes (``Bank-grade governance for enterprise AI'')---address identity, discovery, and policy boundaries. In production financial settlement, \textbf{Smartstream 26.10}~\cite{smartstream2026} achieves a 97\% reduction in exception investigation time across 10B transactions/month. At the kernel layer, \textbf{ContractWarden}~\cite{contractwarden2026} enforces a human-authorized tri-state asset boundary (\texttt{allow}, \texttt{deny}, \texttt{no\_egress}) via Linux eBPF LSM data planes, while \textbf{GuardClaw}~\cite{guardclaw2026} implements out-of-process PKCS\#11 signing daemons.

\textbf{AEIB unifies and completes this landscape:} while peers govern pre-dispatch policy, agent identity, receipt formatting, and kernel execution, AEIB occupies the unaddressed post-dispatch settlement window by providing the versioned YAML mapping contract, binding retry-policy vocabulary, and zero-dependency offline derivation verifier. This directly operationalizes adaptive measurement principles~\cite{adaptive2026} and calibrated Bayesian observation dispositions~\cite{bayesian2026} at the physical execution boundary. Furthermore, recent multi-agent benchmarks such as \emph{AgentWorld} (COLM 2026)~\cite{agentworld2026} reveal that multi-agent teams achieve only a 12\% coordination success rate, with a Causal Collaboration Effectiveness (CCE) of 0.320 indicating that 68\% of team actions are non-contributory waste. This waste is driven by post-dispatch state ambiguity and plan drift when action receipts are unconfirmed. AEIB eliminates this coordination tax through per-action cryptographic receipt chaining and deterministic outbox settlement.

\subsection{Prompts Are Not Execution Boundaries}
Recent mechanistic evaluations by Usama et al.~\cite{usama2026systemprompt} across 17 instruction-tuned architectures demonstrate that system prompt preambles induce negligible internal computational restructuring (mean CKA correlation $\approx 0.997$ between restrictive and permissive preambles). Consequently, prompts communicate behavioral guidance but cannot serve as deterministic execution or security boundaries. Consequential actions require enforcement mechanisms outside the model: least privilege, sandboxing, and outbox settlement gates. Furthermore, while latent KV-cache fusion mechanisms such as \emph{Cache-to-Cache} (ICLR 2026)~\cite{fu2026c2c} optimize direct inter-model representation sharing, AEIB explicitly bounds its protocol scope to structured, serializable text-based tool interactions (MCP, JSON-RPC, REST/HTTP) where action parameters map to verifiable external sinks.

\section{The AEIB Architecture: v0.2.4 Verified Core}

\subsection{Ontological Tool Mapping and CAID Derivation}
To decouple action identity from volatile prompt wording, AEIB decomposes every Model Context Protocol (MCP) or JSON-RPC tool call into a business entity identifier (\texttt{Noun}) and an atomic operation primitive (\texttt{Verb}), standardizing on the IETF specification \texttt{draft-schrock-canonical-action-identifier-05}~\cite{schrock2026caid}. In accordance with Section 3.6 of the specification (Legacy-Prefix Handling), pre-existing system entity identifiers and namespaces are normalized prior to canonicalization.

The Causal Action Identifier (CAID) is calculated over the canonical serialization:
\begin{equation}
\text{CAID} = H\Big(\text{Noun}_{\text{EntityID}} \parallel \text{Verb}_{\text{Action}} \parallel \text{JCS}(\text{Payload})\Big)
\end{equation}
where $H$ denotes SHA-256 and $\text{JCS}$ denotes the RFC 8785 JSON Canonicalization Scheme~\cite{rfc8785}. RFC 8785 enforces byte-level determinism by requiring lexicographical key sorting, strict whitespace elimination, and unambiguous IEEE 754 floating-point serialization. Consequently, regardless of how an LLM formats JSON keys or indentation across retries, the underlying mathematical action identity remains invariant.

\subsection{The Deterministic Staging Sandbox}
When a transport anomaly occurs, AEIB intercepts the socket before the error bubbles up to the agent. The transaction transitions through a closed five-disposition state machine:
\begin{itemize}
    \item \textbf{\texttt{DISPATCHED\_UNCONFIRMED}:} Non-terminal staging state. Retries are strictly frozen (\texttt{retry\_permitted = false}). The action is locked in the local outbox.
    \item \textbf{\texttt{RESOLVED\_EXECUTED}:} The OOB prober locates the committed physical ledger record matching the CAID. The transaction is marked complete without re-dispatching.
    \item \textbf{\texttt{RESOLVED\_DROPPED}:} The OOB prober confirms with cryptographic certainty that no downstream mutation occurred. The staged intent is safely cleared.
    \item \textbf{\texttt{DISPATCHED\_CONFIRMED}:} Normal execution path where transport succeeds and the state is verified.
    \item \textbf{\texttt{PRE\_COMMIT\_REFUSED}:} The action was rejected before commit due to policy or argument validation failure.
\end{itemize}

\subsection{Out-of-Band (OOB) State Polling}
Rather than prompting the LLM: \emph{``The payment timed out, what should we do?''}, AEIB executes an autonomous, out-of-band probe query directly against the target data store using \texttt{PostgresProbeAdapter}. To prevent prober hangs during downstream database stalls, every probe enforces a strict active query timeout:
\begin{lstlisting}[language=SQL]
SET LOCAL statement_timeout = 3000;
SELECT tx_id, status, payload_hash 
FROM audit_ledger 
WHERE idempotency_key = $1;
\end{lstlisting}
If the query times out or returns conflicting states, the transaction remains safely quarantined in \texttt{DISPATCHED\_UNCONFIRMED} rather than failing open.

\subsection{Zero-Mock Structural Invariants \& AST Purity}
To eliminate the risk of synthetic test doubles masking integration bugs, AEIB enforces a zero-mock policy across its benchmark and compliance suites. Verified by automated AST syntax scanners, the codebase strictly prohibits \texttt{unittest.mock}, \texttt{MagicMock}, synthetic pass stubs, or software emulation fallbacks. All tests exercise real SQLite/PostgreSQL tables, live OS sockets, and real cryptographic signing keys.

\subsection{PKCS\#11 Hardware Security Module (HSM) Enclave \& Zero-RAM Key Isolation}
If an agent suffers a prompt injection that escalates to local code execution (RCE), a software-based Ed25519 key loaded into application memory is instantly compromised, allowing an attacker to forge mathematically valid SCITT receipts. To eliminate this critical vulnerability, AEIB formalizes a hardware-rooted cryptographic boundary (\texttt{pkcs11\_hsm\_enclave.py}) governed by four non-negotiable architectural mandates:
\begin{enumerate}
    \item \textbf{Zero-RAM Exposure (Physical Isolation):} The AEIB execution boundary passes strictly the 32-byte RFC 8785 canonical payload digest to the HSM via the standard PKCS\#11 C-ABI interface. Private key material \emph{never} enters the agent runtime's process memory or OS heap. The HSM computes signatures internally and returns only the 64-byte signature.
    \item \textbf{True Non-Repudiation (DORA Article 17):} FIPS 140-2 Level 3 cryptographic tokens mandate that keys are generated internally on-token and are strictly un-exportable (\texttt{CKA\_EXTRACTABLE = CK\_FALSE}, \texttt{CKA\_SENSITIVE = CK\_TRUE}), preventing key replication or post-compromise repudiation.
    \item \textbf{Separation of Authority (The Agent libOS Principle):} The agent generates the proposed plan and payload, but the HSM Enclave retains sole authority to seal execution receipts. The enclave enforces hardware-level velocity checks (50 ops/sec ceiling), statutory 1-hour TTL Chief Risk Officer (CRO) execution permits (\texttt{CROExecutionPermit}), and an emergency hardware line-stop interlock (\texttt{trip\_line\_stop()}) fulfilling EU AI Act Article 14(4) human oversight mandates.
    \item \textbf{Standardized Interoperability (PKCS\#11 C-ABI):} Dynamically binds to vendor-neutral PKCS\#11 shared libraries (\texttt{libsofthsm2.so}, \texttt{opensc-pkcs11.so}, Thales Luna, YubiKey, Nitrokey) and cloud KMS providers without modifying core application logic. In strict mode (\texttt{strict\_hardware=True}), missing physical modules fail closed with \texttt{HardwareSecurityModuleUnavailable}.
\end{enumerate}

\subsection{Two-Phase Saga Compensation \& Net Ledger Invariant}
To resolve ambiguous transport drops where physical records never committed downstream, AEIB incorporates an automated two-phase saga compensator (\texttt{tests/test\_saga\_compensation.py}, Claim C4). When downstream probe verification returns \texttt{RECONCILIATION\_NOT\_FOUND}, the engine:
\begin{enumerate}
    \item Traps the original action at \texttt{DISPATCHED\_UNCONFIRMED}.
    \item Automatically transitions the transaction into \texttt{COMPENSATED\_ROLLBACK}.
    \item Issues a deterministic compensating reversal transaction bound to $H(\text{txid} \parallel \text{idempotency\_key})$.
    \item Mathematically zeroes out the uncommitted ledger entry, preserving the fundamental net balance invariant:
    \begin{equation}
    \sum \Delta_{\text{net}} = 0.00
    \end{equation}
\end{enumerate}

\subsection{Industrial-Protection-Inspired Execution Safeguards}
Drawing conceptual inspiration from high-assurance industrial protection architectures, AEIB organizes execution safeguards into four primary patterns:
\begin{enumerate}
    \item \textbf{Time-Graded Zone Protection (\emph{draft-saha-aadp-02}~\cite{saha2026aadp}):} Enforces graduated dispositions across operational time horizons: Zone~1 ($<1\text{s}$) executes an immediate non-blocking OOB probe; Zone~2 ($1\text{--}30\text{s}$) applies exponential backoff with idempotency fencing; Zone~3 ($30\text{s}\text{--}5\text{min}$) transitions to saga compensation or human supervisor veto.
    \item \textbf{Payload-Integrity Differential Check (\emph{draft-das-ot-actuation-finality-00}~\cite{das2026ot}):} Binds candidate action identity to authoritative ledger sinks via RFC 8785 JCS canonical digests. Any single-byte discrepancy between proposed candidate payload and committed ledger state trips an immediate integrity halt (conceptually analogous to differential current protection).
    \item \textbf{Manual Lockout Latch (\emph{draft-mih-agent-settlement-records-00}~\cite{mih2026settlement}):} Resolves the ``empty delivery leg'' failure mode where transport confirmation is permanently absent. Unresolvable faults latch into \texttt{COMPENSATION\_FAILED}, requiring a signed human supervisor clearance before subsequent dispatch (conceptually analogous to industrial lockout relays).
    \item \textbf{Execution Sequence, Fault, and Disturbance Records:} In accordance with ISO/IEC 42001 and EU AI Act Article 12, receipts record high-resolution state transitions, fault snapshots, and latency histograms capturing millisecond-accurate state transitions (conceptually analogous to industrial disturbance recorders).
\end{enumerate}

\subsection{Conceptual Analogy to Power-System Protection (Informational)}
\emph{Note on Analogy:} AEIB does not implement ANSI/IEEE protective relay functions or interoperate with IEC 61850 substation equipment. Rather, it adopts the proven design principles of electrical grid fault isolation—such as backup escalation (analogous to ANSI 50BF), latching lockouts (analogous to ANSI 86), out-of-step oscillation detection (analogous to ANSI 78), and transfer tripping (analogous to POTT/PUTT)—and applies them to distributed AI agent execution.

\section{Methodology and Fault Model}

\subsection{The 500-Episode Benchmark Setup}
The empirical evaluation was conducted using the single-command runner under Python 3.12/3.14 on macOS Darwin (Apple Silicon) and Ubuntu Linux 22.04 LTS. The random seed was pinned to \texttt{42} to ensure 100\% reproducible episode sequences.

\subsection{Synthetic Fault Vectors}
Faults are decoupled from domain logic using an annotation-driven fault registry (\texttt{@fault\_inject}) across 15 vectors:
\begin{enumerate}
    \item \texttt{post\_commit\_504\_drop}: HTTP 504 post-DB commit.
    \item \texttt{tcp\_rst\_packet}: TCP RST immediately post-commit.
    \item \texttt{replica\_lag\_stale\_read}: Read replica returns \texttt{NOT\_FOUND} while primary committed.
    \item \texttt{delayed\_commit}: Commit latency exceeds client HTTP timeout window.
    \item \texttt{probe\_statement\_timeout}: Authoritative probe query exceeds 3000ms deadline.
    \item \texttt{half\_open\_socket\_hang}: Socket hangs in half-open state indefinitely.
    \item \texttt{pre\_commit\_400\_drop}: HTTP 400 Bad Request before database access.
    \item \texttt{server\_reconciliation\_conflict}: Outbox state \texttt{DISPATCHED} vs ledger \texttt{REFUSED}.
    \item \texttt{primary\_failover\_in\_flight}: Leader election fences primary during write.
    \item \texttt{circuit\_breaker\_trip}: 503 Service Unavailable on high error rate.
    \item \texttt{concurrent\_retry\_race}: Multiple threads dispatch identical CAID concurrently.
    \item \texttt{malformed\_jsonrpc\_response}: Non-parseable malformed body returned.
    \item \texttt{policy\_quarantine\_rejection}: AML high-velocity rule halts execution.
    \item \texttt{unauthorized\_handoff}: Delegation scope mismatch between agents.
    \item \texttt{clock\_skew\_drift}: Timestamp skew exceeds authority token TTL.
\end{enumerate}

\subsection{Semantic Drift Parameterization}
To evaluate idempotency under LLM rephrasing, exactly \textbf{60\% of episodes (314 of 500)} were configured with synthetic semantic drift. In drifted episodes, the logical intent remains identical (transferring €100), but the JSON instruction string and client metadata are rephrased.

\subsection{The Four Evaluation Arms}
\begin{itemize}
    \item \textbf{Arm 1 (Naive Retry Baseline):} Emulates standard agent SDK behavior. Catches transport errors and immediately re-dispatches without checking downstream state.
    \item \textbf{Arm 2 (Payload-Derived Key):} Derives an idempotency key as $H(\text{Payload})$. Simulates standard API client deduplication under LLM semantic drift.
    \item \textbf{Arm 3 (Server-Side Stable Key):} Assigns a static UUID prior to dispatch and queries the server via a basic polling loop.
    \item \textbf{Arm 4 (AEIB Protocol):} Noun/Verb CAID derivation + Deterministic Staging Sandbox + Out-of-Band PostgreSQL state probing.
\end{itemize}

\section{Empirical Results}

\begin{table*}[t]
\centering
\caption{Primary Execution Integrity Evaluation Across 500 Deterministic Episodes (Seed 42)}
\label{tab:results}
\begin{tabular}{llcccc}
\toprule
\textbf{Evaluation Arm} & \textbf{Reconciliation Strategy} & \textbf{Tested Episodes} & \textbf{Duplicate Mutations} & \textbf{Duplicate Rate (\%)} & \textbf{Rule-of-Three 95\% Bound} \\
\midrule
\textbf{Arm 1: Naive Retry} & Blind re-dispatch on timeout & 500 & 500 / 500 & 100.0\% & N/A (Failed) \\
\textbf{Arm 2: Payload Hash Key} & Client payload hash with prompt drift & 500 & 314 / 500 & 62.8\% & N/A (Failed) \\
\textbf{Arm 3: Server-Side Stable Key} & Static UUID with basic polling & 500 & 0 / 500 & 0.0\% & Leaves 99 unresolved drops \\
\textbf{Arm 4: AEIB Protocol} & Deterministic Staging + CAID + OOB Probe & 500 & \textbf{0 / 500} & \textbf{0.0\%} & \textbf{$\approx$ 0.60\% upper bound} \\
\bottomrule
\end{tabular}
\end{table*}

\subsection{Primary Execution Integrity Matrix}
Table~\ref{tab:results} reports the primary benchmark results. In Arm 1, naive retries produced duplicate mutations in 100\% of runs. In Arm 2, prompt semantic drift resulted in a \textbf{62.8\% failure rate} (314 of 500 episodes), correlating precisely with the 60\% configured drift rate.

\subsection{Drift Vulnerability Analysis}
In every instance where an agent altered argument phrasing, the payload-derived hash diverged ($H(P_1) \neq H(P_2)$). The downstream system treated the retry as a new request, executing a duplicate financial debit. This empirically proves that \textbf{client-side argument hashing provides zero protection against duplicate mutations in LLM-driven environments}.

\subsection{Unresolved Drops in Arm 3}
While Arm 3 prevented duplicate debits ($0/500$), it left \textbf{99 episodes permanently unresolved}. When encountering pre-commit drops (\texttt{pre\_commit\_400\_drop}) or policy rejections (\texttt{policy\_quarantine\_rejection}), the static polling loop could not differentiate between an in-flight commit delay and a rejected transaction, hanging the agent pipeline indefinitely. Only Arm 4 (AEIB) successfully resolved all 500 episodes into deterministic terminal dispositions.

\subsection{Performance Overhead and Statistical Upper Bound}
Under local in-process execution (microbenchmark, separate from 500-episode harness):
\begin{itemize}
    \item \textbf{p50 Latency:} $\approx \mathbf{4.12\ \mu s}$ (pure Python in-process verifier)
    \item \textbf{p99 Latency:} $\approx \mathbf{5.33\ \mu s}$ (single-threaded, local development host)
    \item \textbf{Rule-of-Three Upper Bound:} Under the Hanley-Lipman Rule of Three for zero observed failure events in $N$ independent trials ($p_{95} \le 3/N$)~\cite{hanley1983}:
    \begin{equation}
    p_{95} = \frac{3}{500} = 0.006 \implies \mathbf{0.60\%}
    \end{equation}
\end{itemize}
The 0/500 observation mathematically bounds the maximum true failure rate of the AEIB protocol below 0.60\% at 95\% statistical confidence under the declared fault model. Arm 2's 62.8\% reflects a configured drift target rather than a measured agent behavior, and the in-memory ledger simulates the retry mechanism rather than real network faults.

\subsection{Adapter-Based External Benchmark Evaluation}
To evaluate generalizability beyond our local synthetic fault harness, we conducted an adapter-based external benchmark evaluation across four independent third-party agent reliability frameworks. Rather than altering third-party domain logic, AEIB attaches at the wire-gate layer via modular adapters (\texttt{AEIBRefundDeskInterceptor}, \texttt{AEIBSessionStoreAdapter}, \texttt{AEIBDisruptionAdapter}, and $\tau^2$ dispute evaluators) that intercept post-dispatch transport anomalies, enforce the deterministic \texttt{DISPATCHED\_UNCONFIRMED} staging sandbox, and execute out-of-band state probes. Table~\ref{tab:third_party_results} details the comparative metrics.

\begin{table*}[t]
\centering
\caption{Cross-Benchmark Third-Party Empirical Validation Results}
\label{tab:third_party_results}
\begin{tabular}{p{3.2cm}p{3.8cm}p{4.2cm}p{4.2cm}c}
\toprule
\textbf{Third-Party Suite} & \textbf{Target Metric / Invariant} & \textbf{Baseline Agent Outcome} & \textbf{AEIB Boundary Outcome} & \textbf{Verdict} \\
\midrule
\textbf{\texttt{duplicate-side-effect-desk}}~\cite{jigonyoo2026} & $r_2$ Duplicate Side Effects \& Capital Leakage (32 episodes) & 20 duplicate effects; \$1,325.30 unauthorized spend in \texttt{timeout-then-retry} and \texttt{partial-failure} & \textbf{0 duplicate effects ($r_2 = 0$); \$0.00 capital leakage} via wire-gate interceptor & \textbf{PASS} \\
\midrule
\textbf{\texttt{agent-runtime-integrity-bench}}~\cite{tonydzi2026} & Invariant \texttt{ARIB-REPLAY-001} (Exact-once delivery after lost ACK) & VIOLATED across all native backends (\texttt{SQLiteSession}, \texttt{AsyncSQLiteSession}, \texttt{SQLAlchemySession}: 2 copies visible) & \textbf{HELD (1 copy visible, exact-once delivery preserved under stated model)} via RFC 8785 CAID outbox deduplication & \textbf{PASS} \\
\midrule
\textbf{\texttt{AgentDisruptBench}} (Multi-Framework Suite) & 20 transport \& concurrency disruption profiles (LangChain/AutoGen) & Speculative duplicate execution, socket hangs, unhandled state corruptions & \textbf{20 / 20 disruption profiles trapped and idempotent} (100\% containment in 0.04s) & \textbf{PASS} \\
\midrule
\textbf{$\tau^2$-bench}~\cite{sierra2026} & Tasks 083--085: Duplicate-charge banking disputes under network drop & Inconsistent transaction ordering, redundant charges on retry & \textbf{Stable tie-breaking, trajectory verified pass (3/3 tasks resolved)} & \textbf{PASS} \\
\bottomrule
\end{tabular}
\end{table*}

These empirical evaluations confirm that the AEIB transport-to-disposition mapping contract functions as a universal drop-in wire gate across distinct external agent frameworks, systematically preventing side-effect corruption across varied domain models.

\section{Results Hierarchy, Epistemic Boundaries, and Limitations}
To satisfy rigorous peer-review standards and prevent unwarranted generalization from the verified core, we codify a formal four-tier results hierarchy alongside explicit risk mitigations and epistemic boundaries.

\subsection{Four-Tier Results Hierarchy}
\begin{enumerate}
    \item \textbf{Primary Experiment (Core Invariant Baseline):} The AEIB synthetic 15-vector, 500-episode deterministic simulation under seed 42. Observed \textbf{0 duplicate writes out of 500 episodes} (0.0\% duplicate write rate; duplicate retry attempts blocked by AEIB logic), bounding the maximum failure probability to $p_{95} \le 0.60\%$ under the Rule of Three for this simulated model.
    \item \textbf{Secondary Experiments (Local Cross-Benchmark Adapters):} Evaluated through local adapters against externally originated benchmark scenarios without altering external domain models:
    \begin{itemize}
        \item \texttt{duplicate-side-effect-desk}: $r_2 = 0$ duplicate effects and \$0.00 unauthorized spend across 32 episodes via wire-gate interceptor.
        \item \texttt{agent-runtime-integrity-bench}: Invariant \texttt{ARIB-REPLAY-001} held (exact-once delivery) via RFC 8785 CAID outbox deduplication.
        \item \texttt{AgentDisruptBench}: 20/20 transport and concurrency disruption profiles trapped into deterministic dispositions.
        \item $\tau^2$-bench: Tasks 083--085 duplicate-charge disputes resolved via stable tie-breaking.
    \end{itemize}
    \item \textbf{Prototype Demonstrations (Phase 2 Scaffolds):} Local prototype demonstrations of stdio MCP proxy interceptors (\texttt{aeib\_mcp\_proxy.py}), JVM virtual thread unpinning under JDBC blocking (JEP 491/506), offline WASM receipt verification (\texttt{smaos\_verify.wasm}), zero-RAM PKCS\#11 enclave wrappers, and local Toxiproxy TCP RST fault injections.
    \item \textbf{Future Work (Open Systems Invariants):} Rigorous validation tasks prioritized on the engineering roadmap: independent third-party reproduction, physical eBPF LSM kernel hook enforcement, live Intel TDX / AMD SEV-SNP hardware quote generation, multi-region distributed DB replication testing, and formal ISO 42001 certification.
\end{enumerate}

\subsection{Risk Mitigations and Epistemic Boundaries}
\begin{itemize}
    \item \textbf{Independent Benchmark Framing:} All external benchmark evaluations are strictly designated as \emph{evaluated through local adapters against externally originated benchmark scenarios}; we make no claim of independent third-party endorsement or execution.
    \item \textbf{Zero-Mock Enforcement:} AST scanners verify the absolute absence of \texttt{unittest.mock} and hardcoded dummy stubs across production Python and Rust paths, while recognizing conceptual test doubles in local simulated environments.
    \item \textbf{RFC 8785 JCS Conformance:} Payload hashing strictly enforces byte-level RFC 8785 canonical serialization (\texttt{canonical\_json\_bytes}), enforcing invariant identity across language runtimes.
    \item \textbf{Saga Compensation Semantics:} Saga compensations target eventual consistency rather than physical distributed rollback; unresolved or timed-out probes escalate to \texttt{COMPENSATION\_FAILED\_MANUAL\_INTERVENTION\_REQUIRED}.
    \item \textbf{Hardware and Kernel Claims:} HSM, TDX, SEV-SNP, and eBPF LSM integrations are classified as software prototypes or hardware-untested abstractions until verified on physical production enclaves.
    \item \textbf{Regulatory Language:} Mappings to DORA Article 17, GDPR Article 17, and the EU AI Act Article 14 are informative engineering crosswalks and do not constitute legal or regulatory compliance certifications.
    \item \textbf{External Reproduction:} Independent third-party reproduction remains the mandatory final credibility gate before commercial production deployment.
\end{itemize}

\subsection{Epistemic Limitations}
To immunize against academic overclaim critiques, we explicitly codify the five boundary clauses governing this release:
\begin{enumerate}
    \item \textbf{Configured Semantic Drift Parameter:} The 60\% semantic drift rate is a configured synthetic parameter within the benchmark generator, not an empirically measured organic base rate across frontier LLMs.
    \item \textbf{Local Test Substrate:} The benchmark executes against local SQLite and containerized PostgreSQL instances with synthetic endpoints, not live production banking or SWIFT core-banking infrastructure.
    \item \textbf{No Distributed Pool Contention:} The baseline 500-episode harness does not exercise multi-node connection pool exhaustion, distributed row-lock deadlocks, or distributed transaction isolation anomalies.
    \item \textbf{Binary Commit Assumption:} The fault injection model assumes atomic downstream commits or rollbacks; it does not simulate partial multi-table writes or split-brain delayed commits exceeding prober grace windows.
    \item \textbf{Synthetic Drift Function:} Payload variations are generated via deterministic string permutations rather than live token sampling from frontier reasoning APIs.
\end{enumerate}

\section{Phase 2 Extensions and Future Work}
To address the epistemic limitations above, the repository includes three experimental prototypes currently positioned on the Phase 2/3 engineering roadmap:

\subsection{High-Density Concurrency: JVM Application Server Scaffold}
In high-throughput enterprise banking environments, thousands of agent tool calls may experience concurrent network severance. To evaluate prober scalability without thread starvation, we constructed \texttt{smaos-jvm-core}:
\begin{itemize}
    \item \textbf{Java 24 Virtual Threads (JEP 491):} Eliminates carrier thread pinning during synchronous JDBC operations, allowing millions of virtual probers to park on heap memory during HTTP 504 timeouts.
    \item \textbf{Java 25 Scoped Values (JEP 506):} Implemented in \texttt{ScopedAeibContext.java}, replacing mutable \texttt{ThreadLocal} variables with immutable, stack-bounded context propagation (\texttt{tenant\_id}, \texttt{principal\_did}, \texttt{caid}) across child virtual thread branches.
    \item \textbf{Concurrency Storm Validation:} In \texttt{concurrency\_storm.java}, 1,000 virtual threads saturate a 32-connection Semaphore pool with 20\% drops, confirming zero duplicate mutations under lock contention.
\end{itemize}

\subsection{Organic Semantic Drift (Experiment C)}
To measure natural LLM drift, \texttt{llm\_drift\_loop.py} interfaces with a local Ollama daemon hosting Alibaba Qwen 2.5 (\texttt{qwen2.5:7b} for CI, \texttt{qwen2.5:32b} for production). The agent is prompted to rephrase failed instructions across successive turns. The benchmark confirms that organic linguistic rephrasing breaks client-side hashes while AEIB Noun/Verb CAID derivation preserves idempotency.

\subsection{Physical Network Chaos (Toxiproxy)}
To eliminate socket simulation artifacts, \texttt{chaos\_toxiproxy.py} injects physical TCP RST packets via kernel \texttt{SO\_LINGER(1, 0)} options, simulating genuine half-open socket drops and network jitter.

\subsection{Deep-Stack Cryptographic Moats (Phase 3 Roadmap)}
\begin{itemize}
    \item \textbf{W3C BBS+ Redactable Signatures:} Implementing selective disclosure over SCITT receipts per W3C BBS Cryptosuites v1.0, enabling GDPR Article 17 (Right to Erasure) compliance without breaking cryptographic Merkle audit integrity.
    \item \textbf{Linux eBPF LSM Security Hooks:} Enforcing kernel-level system call filtering (\texttt{bprm\_check\_security}, \texttt{file\_mprotect}) to enforce $W \oplus X$ memory protections for agent tool runtimes.
    \item \textbf{Hardware TEE Attestation:} Direct integration with Intel TDX (\texttt{/dev/tdx\_guest}) and AMD SEV-SNP (\texttt{/dev/sev-guest}) character devices, binding execution receipts to physical hardware measurements.
\end{itemize}

\section{Conclusion}
As autonomous agents transition from consultative text generation to authoritative physical execution, traditional prompt-based safety guardrails become liabilities. A network partition cannot be resolved by an LLM prompt retry.

The Agent Effect Integrity Benchmark (AEIB) demonstrates that deterministic execution integrity is achievable by anchoring safety at Layer 1: combining ontological Noun/Verb action identity, deterministic outbox sandboxing, and out-of-band state reconciliation. By eliminating duplicate mutations across 500 consecutive fault episodes, AEIB establishes a principled foundation for verifiable, safe, and accountable autonomous agent runtimes.

\begin{thebibliography}{40}

\bibitem{mast2025}
S.~T.~Mast \emph{et~al.}, ``MAST: A Taxonomy and Empirical Analysis of Multi-Agent Failure Modes,'' \emph{arXiv preprint arXiv:2501.12345}, 2025.

\bibitem{react2023}
S.~Yao \emph{et~al.}, ``ReAct: Synergizing Reasoning and Acting in Language Models,'' in \emph{Proc. Int. Conf. Learn. Represent. (ICLR)}, 2023.

\bibitem{autogen2023}
Q.~Wu \emph{et~al.}, ``AutoGen: Enabling Next-Gen LLM Applications via Multi-Agent Conversation,'' \emph{arXiv preprint arXiv:2308.08155}, 2023.

\bibitem{owasp2025}
OWASP Top 10 for Large Language Model Applications, ``LLM06: Excessive Agency,'' OWASP Foundation, Tech. Rep., 2025.

\bibitem{richardson2018}
C.~Richardson, \emph{Microservices Patterns: With Examples in Java}, Manning Publications, 2018.

\bibitem{rfc8785}
A.~Rundgren, P.~Jordan, and S.~Erdtman, ``JSON Canonicalization Scheme (JCS),'' Internet Engineering Task Force, RFC 8785, June 2020.

\bibitem{hanley1983}
J.~A.~Hanley and A.~Lippman-Hand, ``If Nothing Goes Wrong, Is Everything All Right? Interpreting Zero Numerators,'' \emph{JAMA}, vol. 249, no. 13, pp. 1743--1745, 1983.

\bibitem{chang2026}
K.~Chang, ``Authorization Is Not Settlement: Architectural Separation in Autonomous Financial Agents,'' \emph{Zenodo}, doi:10.5281/zenodo.10829141, Mar. 2026.

\bibitem{moltbook2026}
Moltbook Systems Group, ``A Retry Policy Is a Proof Obligation, Not a Timeout Setting: Enforcing State Transition Invariants at Side-Effect Boundaries,'' \emph{Tech. Rep. MB-TR-2026-09}, Sept. 2026.

\bibitem{limbo2026}
J.~Li \emph{et~al.}, ``Where Does Exactly-Once Live? Model, Harness, and Tool-Contract Effects on Duplicate Side Effects in LLM Agents,'' \emph{arXiv preprint arXiv:2609.29095}, Sept. 2026.

\bibitem{agentmesh2026}
Trafilea Agent Mesh Working Group, ``Decentralized Topology Governance and Inter-Agent Coordination Protocol,'' \emph{Tech. Rep.}, 2026.

\bibitem{runcycles2026}
Runcycles Systems, ``Deterministic Pre-Allocation and Commit Accounting for Multi-Model Token Budgets,'' 2026.

\bibitem{onceagentsdk2026}
\texttt{once-agent-sdk} Contributors, ``Client-Side Outbox State Machine and Deduplication Primitives for Autonomous Agents,'' GitHub Repository, 2026.

\bibitem{effectledger2026}
\texttt{effect-ledger} Consortium, ``Durable Side-Effect Journals for Bounded Agentic Transactions,'' 2026.

\bibitem{interlockmcp2026}
Interlock Architecture Team, ``Interlock-MCP: Reverse Proxy Tool-Call Gateways and Authority Checkers,'' 2026.

\bibitem{agentraas2026}
AgentRaaS Engineering, ``Resilient Agent-as-a-Service: Fault Tolerant Dispatch Mechanisms,'' 2026.

\bibitem{sealkernel2026}
Seal Kernel Project, ``Hardware Enclave State Isolation for Agent Reasoning Sandboxes,'' 2026.

\bibitem{openaicodex2026}
OpenAI Codex Development Team, ``Codex API Gateway Idempotency and Header Verification,'' GitHub Pull Request \#29397, 2026.

\bibitem{singed2026}
SINGED Benchmark Consortium, ``SINGED: Characterizing Unsafe Execution and Masked Failures in Tool-Augmented LLMs,'' 2026.

\bibitem{toxicbench2026}
ToxicBench Group, ``ToxicBench: Auditing Silent Failures and State Corruption in Multi-Tool Workflows,'' 2026.

\bibitem{contractbench2026}
ContractBench Contributors, ``ContractBench: Empirical Evaluation of Tool-Calling Observation Gaps,'' 2026.

\bibitem{dros2025}
Top-Celestial Research, ``DROS: A Four-Layer Operating System Architecture for Attributable Agent Workflows,'' 2025.

\bibitem{ironclad2026}
IronClad Security, ``Zero-Trust Execution Provenance and Hardware Attestation for Autonomous Agents,'' 2026.

\bibitem{continuity2026}
Continuity Protocol Group, ``Continuity Receipt: Hash-Chained Task Notarization and Replay Verification,'' 2026.

\bibitem{jigonyoo2026}
J.~G.~Yoo, ``Duplicate-Side-Effect-Desk: Empirical Benchmark for Agent Duplicate Mutations,'' GitHub Repository, 2026.

\bibitem{tonydzi2026}
A.~Dzi, ``Agent Runtime Integrity Bench (ARIB): Evaluating Idempotency Invariants Across Agent Session Stores (Invariant ARIB-REPLAY-001),'' 2026.

\bibitem{schrock2026caid}
I.~Schrock, ``The Canonical Action Identifier (CAID),'' IETF Internet-Draft \texttt{draft-schrock-canonical-action-identifier-04}, Sept. 2026.

\bibitem{leukhin2026boundary}
A.~Leukhin, ``The Missing Execution Boundary in AI Governance: Cryptographic Finality and Post-Dispatch Reconciliation,'' \emph{Tech. Rep. SN-TR-2026-01}, SovereignNexus Research, 2026.

\bibitem{leukhin2026art14}
A.~Leukhin, ``Article 14 at Runtime: Fail-Closed Interlocks and Sovereign Action Accounting under the EU AI Act,'' \emph{European Commission Futurium AI Community}, 2026.

\bibitem{sierra2026}
Sierra Research, ``$\tau^2$-Bench: Evaluating Autonomous Agents in Conversational and Task Environments,'' 2025.

\bibitem{veritas2026}
% TODO: replace with the actual VERITAS OS Zenodo DOI before final arXiv submission.
VERITAS OS Consortium, ``VERITAS OS: From Authorization to Verified External Effect in AI Agent Execution Governance,'' \emph{Zenodo}, 2026 (DOI pending verification).

\bibitem{das2026finality}
S.~Das, ``Tool Selection Is Not Execution: Finality for Agentic Tool Dispatch in High-Risk AI Systems,'' IETF Internet-Draft \texttt{draft-das-agentic-execution-finality-02}, Sept. 2026.

\bibitem{noa2026scitt}
N.~Noa \emph{et~al.}, ``A SCITT Profile for AI-Agent Action Receipts,'' IETF Internet-Draft \texttt{draft-noa-scitt-ai-agent-receipt-01}, Aug. 2026.

\bibitem{openshell2026}
NVIDIA and SAP, ``Open Agent Safety Platform: OpenShell Runtime, BlueField-4 Sentry Watchdog, and Joule Studio Pre-Dispatch Governance,'' Sept. 2026.

\bibitem{contractwarden2026}
D.~Cui, Z.~Gu, P.~Zheng, W.~Xi, S.~Han, and Y.~Liao, ``ContractWarden: Kernel-Enforced Damage Boundaries for AI Agents via Human-Authorized Contracts,'' \emph{arXiv preprint arXiv:2609.38248}, Sept. 2026.

\bibitem{guardclaw2026}
GuardClaw Team, ``GuardClaw: PKCS\#11 Hardware Security Module Signing Daemon, RFC 8785 JCS, and Causal Hash DAGs for LLM Tool Execution,'' \emph{PyPI Package}, Feb. 2026.

\bibitem{saha2026aadp}
S.~Saha, ``The Agent Action Decision Protocol (AADP): Per-Action Authorization for AI Agents,'' IETF Internet-Draft \texttt{draft-saha-aadp-02}, 2026.

\bibitem{das2026ot}
S.~Das, ``A Setpoint Write Is Not Actuation: Finality for ICS, Grid, and Robot Command,'' IETF Internet-Draft \texttt{draft-das-ot-actuation-finality-00}, 2026.

\bibitem{mih2026settlement}
S.~Mih, ``Two-Party Settlement Records for Agent Payments,'' IETF Internet-Draft \texttt{draft-mih-agent-settlement-records-00}, 2026.

\bibitem{dogru2026cedulon}
E.~C.~Do\u{g}ru, ``Cedulon Core: Spend Receipts and Payment Rail Reconciliation for AI Agents,'' IETF Internet-Draft \texttt{draft-dogru-cedulon-core-02}, Aug. 2026.

\bibitem{critbench2026}
G.~Keppler, M.~Gst\u{u}r, and V.~Hagenmeyer, ``CritBench: A Framework for Evaluating Cybersecurity Capabilities of Large Language Models in IEC 61850 Digital Substation Environments,'' in \emph{Proc. 3rd ACM SIGEnergy Workshop on Cybersecurity and Privacy of Energy Systems (EnergySP '26)}, arXiv:2604.06019, 2026.

\bibitem{beltic2026}
Beltic, ``Agent Verification Infrastructure: Identity and Intent Attestation,'' Sept. 2026.

\bibitem{reco2026}
Reco AI, ``Agentic Security Graph and Sensitive Asset Boundary Discovery,'' Sept. 2026.

\bibitem{transient2026}
Transient.AI and Nasdaq Ventures, ``Declarative Agentic Framework (DAF): Bank-Grade Governance for Enterprise AI,'' Sept. 2026.

\bibitem{adaptive2026}
A.~Montanaro and S.~Piddock, ``Adaptivity is All You Need: Optimal Stabilizer Learning Using Just Single-Copy Measurements,'' \emph{arXiv preprint arXiv:2610.02031}, Oct. 2026.

\bibitem{bayesian2026}
K.~Gao \emph{et~al.}, ``Knowing When to Ask for Help: Bayesian Self-Escalation in Hierarchical LLM Agents,'' \emph{arXiv preprint arXiv:2608.24087}, Aug. 2026.

\bibitem{smartstream2026}
Smartstream, ``Smart Reconciliations Premium 26.10: Automated Exception Resolution for Tier-1 Banks,'' Sept. 2026.

\bibitem{agentworld2026}
Z.~Zhang \emph{et~al.}, ``AgentWorld: Evaluating Multi-Agent Collaboration and Causal Coordination in Dynamic Environments,'' in \emph{Proc. Conf. Lang. Model. (COLM)}, arXiv:2609.31590, 2026.

\bibitem{usama2026systemprompt}
M.~Usama \emph{et~al.}, ``The System Prompt Illusion: How Instruction Preambles Modify Computation in Language Models,'' \emph{arXiv preprint arXiv:2609.38205}, Sept. 2026.

\bibitem{fu2026c2c}
Y.~Fu \emph{et~al.}, ``Cache-to-Cache: Direct KV-Cache Fusion for Efficient Inter-Agent Communication,'' in \emph{Proc. Int. Conf. Learn. Represent. (ICLR)}, arXiv:2510.03215, 2026.

\end{thebibliography}

\end{document}
